\documentclass[11pt,a4paper]{article}

\usepackage[margin=1in]{geometry}
\usepackage{graphicx}
\usepackage{float}
\usepackage{wrapfig}
\usepackage{caption}
\usepackage{subcaption}
\usepackage{booktabs}
\usepackage{array}
\usepackage{enumitem}
\usepackage[table]{xcolor}
\usepackage{amsmath, amssymb}
\usepackage{listings}
\usepackage{algorithm}
\usepackage{algpseudocode}
\usepackage{comment}
\usepackage[numbers]{natbib}
\usepackage{hyperref}

\definecolor{customblue}{RGB}{0, 76, 153}
\definecolor{softgreen}{RGB}{226, 245, 236}
\definecolor{softred}{RGB}{255, 232, 232}
\definecolor{softyellow}{RGB}{255, 247, 214}
\DeclareMathOperator*{\argmax}{arg\,max}

\lstset{
    basicstyle=\ttfamily\small,
    keywordstyle=\color{customblue}\bfseries,
    commentstyle=\color{teal!70!black},
    stringstyle=\color{red!70!black},
    frame=single,
    breaklines=true,
    showstringspaces=false,
    columns=fullflexible
}

\hypersetup{
    colorlinks=true,
    linkcolor=customblue,
    citecolor=customblue,
    urlcolor=teal
}

\title{CSE200 LaTeX Tutorial II}
\author{Mohammad Sadat Hossain}
\date{June 2026}

\begin{document}

\maketitle

\section{Text Formatting}

Package note: bold, underline, and emphasis commands are built in. Text color
needs the \texttt{xcolor} package. Colored table rows and cells need the
\texttt{table} option for \texttt{xcolor}, as shown in the preamble.

The commands below are worth knowing for text formatting.

\begin{itemize}
    \item \verb|\textbf{...}| gives \textbf{bold text}.
    \item \verb|\underline{...}| gives \underline{underlined text}.
    \item \verb|\emph{...}| gives \emph{emphasized text}; its appearance may
    change depending on context.
\end{itemize}

Example sentence: A technical report should use \textbf{important terms},
\underline{precise labels}, and \emph{consistent emphasis} carefully.

Use \texttt{xcolor} when color has a clear purpose: emphasis, status, or
readability. This sentence has \textcolor{red}{red text},
\textcolor{customblue}{a custom phrase}, and a
\colorbox{softyellow}{short highlighted note}.

\begin{table}[H]
    \centering
    \rowcolors{2}{gray!10}{white}
    \begin{tabular}{lll}
        \toprule
        Item & Result & Comment \\
        \midrule
        Figure reference & \cellcolor{softgreen}Correct & Label placed after caption \\
        Image size & \cellcolor{softred}Needs work & Width too large \\
        Citation & \cellcolor{softgreen}Correct & Source appears in bibliography \\
        \bottomrule
    \end{tabular}
    \caption{Color used to communicate status in a table.}
    \label{tab:color-status}
\end{table}

\section{Nested and Mixed Lists}

Package note: basic \texttt{itemize} and \texttt{enumerate} lists are built in.
Custom list labels in this section use the \texttt{enumitem} package.

Technical documents often combine numbered lists, bullets, custom symbols, and
lettered sublists. The \texttt{enumitem} package gives reliable control over
list labels.

\begin{enumerate}[label=\arabic*.]
    \item System Setup
    \begin{itemize}
        \item Install software packages
        \item Configure hardware
        \item[!] Verify connections
    \end{itemize}

    \item Simulation Steps
    \begin{enumerate}[label=(\alph*)]
        \item Initialize parameters
        \item Run simulation
        \begin{itemize}[label=-]
            \item Monitor CPU usage
            \item[$\rightarrow$] Record intermediate results
        \end{itemize}
        \item Analyze outputs
    \end{enumerate}

    \item Reporting
    \begin{itemize}
        \item Prepare charts
        \item[\#] Document observations
        \item Summarize findings
        \begin{enumerate}[label=(\alph*)]
            \item Key insights
            \item Recommendations
        \end{enumerate}
    \end{itemize}
\end{enumerate}

\section{Image Files and the \texttt{figure} Float}

Package note: image insertion needs \verb|\usepackage{graphicx}|. The
\texttt{figure} environment itself is built in.

Use \verb|\includegraphics| to insert an image. Put it inside a
\texttt{figure} environment when it needs a number, caption, label, and
cross-reference.

\begin{figure}[ht]
    \centering
    \includegraphics[width=0.55\linewidth]{images/bangladesh_satellite.jpg}
    \caption{Bangladesh observed by NASA's Terra satellite.}
    \label{fig:bangladesh}
\end{figure}

Figure~\ref{fig:bangladesh} is referenced from the paragraph. The label comes
after the caption so that LaTeX stores the correct figure number.

\subsection{Image Size, Scale, and Rotation}

\begin{figure}[H]
    \centering
    \includegraphics[width=0.28\linewidth]{images/padma_river.jpg}
    \hfill
    \includegraphics[scale=0.16]{images/hubble_cluster.jpg}
    \hfill
    \includegraphics[width=0.28\linewidth, angle=8]{images/smoke_satellite.png}
    \caption{The same figure can demonstrate width, scale, and angle options.}
    \label{fig:size-scale-angle}
\end{figure}

Prefer \verb|width=...| with \verb|\linewidth| for predictable layouts. Use
\verb|scale=...| when you want to resize by a factor, and \verb|angle=...| for
rotation.

\section{Float Placement: \texttt{h}, \texttt{t}, \texttt{b}, \texttt{p},
\texttt{!}, and \texttt{H}}

LaTeX floats are allowed to move. The placement option is a suggestion, not a
command, except for \verb|[H]| from the \texttt{float} package.
Package note: \texttt{h}, \texttt{t}, \texttt{b}, \texttt{p}, and \texttt{!}
are built in; \texttt{H} requires \verb|\usepackage{float}|.

\begin{table}[H]
    \centering
    \begin{tabular}{clp{0.48\linewidth}}
        \toprule
        Option & Meaning & Teaching note \\
        \midrule
        \texttt{h} & here & Try to place the float near this point. \\
        \texttt{t} & top & Allow placement at the top of a page. \\
        \texttt{b} & bottom & Allow placement at the bottom of a page. \\
        \texttt{p} & page & Allow a separate page containing only floats. \\
        \texttt{!} & override & Relax some internal restrictions such as float
        area limits. It does not mean ``force exactly here.'' \\
        \texttt{H} & exactly here & Requires \texttt{float}; useful for demos
        and draft documents, but should not be overused in polished writing. \\
        \bottomrule
    \end{tabular}
    \caption{Common float placement options.}
    \label{tab:float-options}
\end{table}

\subsection{Placement Examples}

\begin{figure}[t]
    \centering
    \includegraphics[width=0.38\linewidth]{images/hubble_cluster.jpg}
    \caption{A top-position float using \texttt{[t]}.}
    \label{fig:top-float}
\end{figure}

\begin{figure}[b]
    \centering
    \includegraphics[width=0.38\linewidth]{images/padma_river.jpg}
    \caption{A bottom-position float using \texttt{[b]}.}
    \label{fig:bottom-float}
\end{figure}

\begin{figure}[!htbp]
    \centering
    \includegraphics[width=0.45\linewidth]{images/smoke_satellite.png}
    \caption{A flexible float using \texttt{[!htbp]}.}
    \label{fig:flexible-float}
\end{figure}

There are no standard capital variants such as \texttt{T} or \texttt{B} for
ordinary article floats. The important capital option for this lesson is
\verb|H|, provided by the \texttt{float} package.

\section{Wrapped Figures}

Package note: wrapped figures need \verb|\usepackage{wrapfig}|.

\begin{wrapfigure}{r}{0.36\linewidth}
    \centering
    \includegraphics[width=\linewidth]{images/padma_river.jpg}
    \caption{Wrapped image.}
    \label{fig:wrapped-padma}
\end{wrapfigure}

A wrapped figure lets text flow beside an image. This is useful when a small
figure belongs to the same paragraph as its explanation. Keep enough text
after the wrapped figure; otherwise, the next paragraph may begin too close to
the float. Also avoid making wrapped images too wide. A width between
\verb|0.25\linewidth| and \verb|0.40\linewidth| is often a good starting point.
Figure~\ref{fig:wrapped-padma} shows a compact layout that is useful when an
image and its explanation should stay visually close.

\clearpage

\section{Subfigures}

Package note: subfigures in this demo need \verb|\usepackage{subcaption}|.

Subfigures need the \texttt{subcaption} package. Each image sits inside a
\texttt{subfigure} block. The block width controls the slot; inside the slot,
use \verb|width=\linewidth| so the image fills that slot.

\begin{figure}[H]
    \centering
    \begin{subfigure}{0.45\linewidth}
        \centering
        \includegraphics[width=\linewidth]{images/bangladesh_satellite.jpg}
        \caption{Country scale}
        \label{fig:sub-country}
    \end{subfigure}
    \hfill
    \begin{subfigure}{0.45\linewidth}
        \centering
        \includegraphics[width=\linewidth]{images/padma_river.jpg}
        \caption{River scale}
        \label{fig:sub-river}
    \end{subfigure}
    \vspace{0.35cm}
    \begin{subfigure}{0.45\linewidth}
        \centering
        \includegraphics[width=\linewidth]{images/hubble_cluster.jpg}
        \caption{Space scale}   
        \label{fig:sub-space}
    \end{subfigure}
    \hfill
    \begin{subfigure}{0.45\linewidth}
        \centering
        \includegraphics[width=\linewidth]{images/smoke_satellite.png}
        \caption{Atmospheric view}
        \label{fig:sub-smoke}
    \end{subfigure}
    \caption{Four subfigures arranged as a two-by-two grid.}
    \label{fig:four-subfigures}
\end{figure}

\textbf{\textcolor{red}{WARNING: DO NOT put blank line between subfigures. It starts a new paragraph, and forces the next subfigure onto a new line, no matter the sizes.}}

We can reference the full figure as Figure~\ref{fig:four-subfigures}, or a
single part such as Figure~\ref{fig:sub-river}. The main layout trick is to
make each subfigure slightly less than half the line width. The \verb|\hfill|
command separates the two columns, and the blank line with \verb|\vspace|
starts the second row.

\section{Hyperlinks}

Package note: links and clickable references need \verb|\usepackage{hyperref}|.

The \texttt{hyperref} package makes cross-references, citations, and web links
clickable. Use \verb|\href| when the displayed text should be different from
the URL, and \verb|\url| when the URL itself should be printed.

Example: \href{https://www.overleaf.com/learn/latex/Main_Page}{Overleaf Learn
LaTeX}. A raw URL can also be printed with \verb|\url|:
\url{https://www.overleaf.com}.

\section{Footnotes}

Package note: ordinary footnotes are built in; no extra package is needed.

Use \verb|\footnote{...}| for short side comments, source notes, or small
clarifications that should not interrupt the main paragraph.\footnote{This is a
regular footnote. It appears at the bottom of the page.} Footnotes are numbered
automatically, and LaTeX handles their placement.

\begin{verbatim}
This claim needs a short explanation.\footnote{Write the note here.}
\end{verbatim}

If the same note must be referenced twice, one simple approach is to store the
footnote number and reuse it:

\begin{verbatim}
First mention.\footnote{\label{fn:tool}A repeated note.}
Second mention.\footnotemark[\ref{fn:tool}]
\end{verbatim}

Avoid long footnotes in short reports. If the note becomes a full paragraph,
it probably belongs in the main text.

\section{Paper Citations and Bibliography}

Package note: this compiled demo uses BibTeX together with the
\texttt{natbib} package. Plain BibTeX can work without a citation package, but
\texttt{natbib} adds better citation commands while still using a
\texttt{.bib} file and BibTeX.

For CSE reports, students must know how to cite papers, not just websites. The
workflow is: keep source details in a \texttt{.bib} database, cite those keys
in the text, and print the bibliography at the end of the document.

\subsection{Plain BibTeX Starting Point}

Plain BibTeX begins with the built-in \verb|\cite{...}| command. These lines
are intentionally shown as comments because the compiled demo below uses the
more expressive \texttt{natbib} commands instead.

\begin{verbatim}
% Single source with plain BibTeX:
% \cite{vaswani2017attention}

% Multiple sources with plain BibTeX:
% \cite{krizhevsky2012imagenet,he2016deep}
\end{verbatim}

\subsection{Natbib Commands in Action}

Use \verb|\citet| when the author name is part of the sentence.

\begin{verbatim}
\citet{vaswani2017attention} introduced the Transformer architecture.
\end{verbatim}

Compiled result: \citet{vaswani2017attention} introduced the Transformer
architecture.

Use \verb|\citep| when the citation is parenthetical evidence at the end of a
claim.

\begin{verbatim}
Residual learning became a standard CNN technique \citep{he2016deep}.
\end{verbatim}

Compiled result: residual learning became a standard CNN technique
\citep{he2016deep}.

Multiple papers can still be cited together by putting several keys inside the
same command.

\begin{verbatim}
\citep{krizhevsky2012imagenet,he2016deep}
\end{verbatim}

Compiled result: image classification and residual networks can be cited
together \citep{krizhevsky2012imagenet,he2016deep}.

Sometimes only the author name or only the year is needed.

\begin{verbatim}
\citeauthor{he2016deep} proposed residual learning in \citeyear{he2016deep}.
\end{verbatim}

Compiled result: \citeauthor{he2016deep} proposed residual learning in
\citeyear{he2016deep}.

The package option controls the broad citation style. This document uses the
numeric option:

\begin{verbatim}
\usepackage[numbers]{natbib}
\end{verbatim}

For an author-year demo, one can try the author-year option together with a
compatible bibliography style such as \texttt{plainnat}.

\begin{verbatim}
\usepackage[authoryear,round]{natbib}
\end{verbatim}

\subsection{Including Uncited Sources}

When a source should appear in the bibliography even though it is not cited in
the paragraph, use \verb|\nocite{...}|. This command does not print anything at
the current location; it only asks LaTeX and BibTeX to include that entry in the
reference list.

\begin{verbatim}
\nocite{lamport1994latex}
\end{verbatim}

\nocite{lamport1994latex}

The command \verb|\nocite{*}| prints every entry from the \texttt{.bib} file.
That can be useful for a bibliography demo, but it is usually too broad for a
finished report.

\subsection{Bibliography Files}

The source details live in \texttt{references.bib}. A paper entry usually
contains fields such as \texttt{author}, \texttt{title}, \texttt{booktitle} or
\texttt{journal}, \texttt{year}, and sometimes \texttt{pages}, \texttt{doi}, or
\texttt{url}.

\begin{verbatim}
@inproceedings{vaswani2017attention,
  author = {Ashish Vaswani and others},
  title = {Attention Is All You Need},
  booktitle = {Advances in Neural Information Processing Systems},
  year = {2017}
}
\end{verbatim}

\subsection{Styles and Compile Order}

Change the bibliography style to change how references are printed. This
document currently uses:

\begin{verbatim}
\bibliographystyle{plainnat}
\bibliography{references}
\end{verbatim}

The style command selects the formatting style. The bibliography command tells
BibTeX to read \texttt{references.bib}.
Common BibTeX-style options include:

\begin{itemize}
    \item \texttt{plainnat}: natbib-friendly plain numeric or author-year style.
    \item \texttt{plain}: numeric labels, sorted alphabetically by author.
    \item \texttt{abbrv}: like \texttt{plain}, but with abbreviated names.
    \item \texttt{unsrt}: numeric labels ordered by first citation.
    \item \texttt{ieeetr}: IEEE-like numeric style, common in engineering.
    \item \texttt{apalike}: author-year-like references for simple reports.
\end{itemize}

\section{Algorithms, Pseudocode, and Code Listings}

Package note: pseudocode here uses the \texttt{algorithm} and
\texttt{algpseudocode} packages. Code listings use the \texttt{listings}
package.

Technical reports often need pseudocode or small code snippets, and Overleaf
also treats code listing as a standard topic. Use \texttt{listings} for code
because it compiles normally. Avoid \texttt{minted} in an introductory setting
unless shell escape is explicitly allowed.

\begin{algorithm}[H]
    \caption{Binary Search}
    \label{alg:binary-search}
    \begin{algorithmic}[1]
        \Require Sorted array $A$, target value $x$
        \Ensure Index of $x$, or $-1$ if not found
        \State $low \gets 1$, $high \gets n$
        \While{$low \leq high$}
            \State $mid \gets \lfloor (low + high)/2 \rfloor$
            \If{$A[mid] = x$}
                \State \Return $mid$
            \ElsIf{$A[mid] < x$}
                \State $low \gets mid + 1$
            \Else
                \State $high \gets mid - 1$
            \EndIf
        \EndWhile
        \State \Return $-1$
    \end{algorithmic}
\end{algorithm}

\begin{lstlisting}[language=Python, caption={A small Python function.}, label={lst:python-square}]
def greet_user(name):
    # Create a greeting message using an f-string
    message = f"Hello, {name}! Welcome to Python."

    # Print the greeting message
    print(message)

    # Return the message in case it's needed elsewhere
    return message


# Example usage
user_name = "Alice"
greet_user(user_name)
\end{lstlisting}

Algorithm~\ref{alg:binary-search} and Listing~\ref{lst:python-square} can be
referenced just like figures and tables.

\section{Two-Column Documents}

Package note: document-wide two-column layout can be requested in the document
class option, e.g., \verb|\documentclass[twocolumn]{article}|. The separate
demo also uses \texttt{graphicx}, \texttt{subcaption}, \texttt{booktabs}, and
\texttt{hyperref}. Open and compile \texttt{two-column-demo.tex}; it uses a
full two-column article with \verb|\documentclass[twocolumn]{article}|, normal
figures, and a wide figure spanning both columns with \texttt{figure*}.

\section{Debugging Checklist for This Lesson}

\begin{itemize}
    \item If a reference shows as \texttt{??}, compile LaTeX again.
    \item If a citation shows as \texttt{?}, check the citation.
    \item If an image is missing, check the file name, extension, and folder.
    \item If a float is far away, try \verb|[!htbp]| before forcing \verb|[H]|.
    \item If a subfigure overflows, reduce each subfigure width before changing
    manual spacing.
    \item If colored table cells fail, load \verb|\usepackage[table]{xcolor}|.
\end{itemize}

\bibliographystyle{plainnat}
\bibliography{references}

\end{document}
